\section{Conclusão}

O trabalho implementou e comparou o GRASP clássico e o reativo para o CVRP. A calibração selecionou $\alpha=0.25$. Nos experimentos finais, o reativo obteve gaps agregados ligeiramente menores, mas executou cerca de 43\% mais lentamente. Os gráficos não mostraram dominância: o clássico obteve mais sucessos em alcançar o ótimo em A-n33-k5 e B-n31-k5.

Como trabalhos futuros, recomenda-se ampliar as repetições e testar vizinhanças maiores ou path-relinking.
